\section{Fractional-moment decay and sampling-law separation}
\label{sec:fractional}

\subsection{The pair inequality and admissible moment tests}

\begin{lemma}[Sharp fractional pair inequality]\label{lem:pair}
For $x,y>0$ and $0<p<1$,
\begin{equation}
 (x+y)\bigl[x^p+y^p-(x+y)^p\bigr]
 \geq(2-2^p)(xy^p+yx^p).
 \label{eq:pair}
\end{equation}
Equality holds exactly when $x=y$.
\end{lemma}

\begin{proof}
Put $u=x/(x+y)$, $v=1-u$, and $A=2^p-1\in(0,1)$.
After division by $(x+y)^{p+1}$, \eqref{eq:pair} is equivalent to
$g(u)\geq1$, where
\[
 g(u)=A(u^p+v^p)+(1-A)(u^{p+1}+v^{p+1}).
\]
By symmetry it suffices to consider $0<u<1/2$. The continuous endpoint
values are $g(0)=g(1/2)=1$, whereas $g'(0+)=+\infty$ and $g'(1/2)=0$.
The sign of $g''(u)$ is the sign of
\[
 (1-A)(p+1)-A(1-p)R(u),\qquad
 R(u)=\frac{u^{p-2}+v^{p-2}}{u^{p-1}+v^{p-1}}
     =\frac{u^{2-p}+v^{2-p}}{uv(u^{1-p}+v^{1-p})}.
\]
On $(0,1/2)$ the numerator in the final expression strictly decreases,
and each positive factor in its denominator strictly increases. Hence $R$
strictly decreases and $g''$ changes sign at most once, from negative to
positive. This change must occur. Since $R(u)\to+\infty$ as
$u\downarrow0$, $g''<0$ near $0$; equivalently, $g'(0+)=+\infty$ is
incompatible with $g'$ increasing on $(0,1/2)$ to the value $g'(1/2)=0$.
Nor can $g''$ remain negative: then $g'$ would decrease to zero while
remaining positive, contradicting the equal endpoint values of $g$.
On the final interval where $g''>0$, the derivative approaches zero from
below. On the initial interval it therefore crosses zero exactly once.
Thus $g$ first increases and then decreases to its endpoint value $1$,
and $g(u)>1$ in $(0,1/2)$. This proves the inequality and its equality case.
\end{proof}

Estimates of the subadditivity defect $x^p+y^p-(x+y)^p$ are standard
tools in the fractional-moment analysis of Smoluchowski equations; see
\citet{EscobedoMischlerPerthame2002},
\citet{FournierLaurencot2005,FournierLaurencot2006}, and
\citet{EscobedoMischlerRodriguezRicard2005}. Inequality \eqref{eq:pair}
is a weighted variant of these estimates, with the additive weight $x+y$
built in. The contribution here is the explicit constant $2-2^p$ with its
equality case, and its use below as a uniform coefficient for
parent-dependent daughters.

\begin{lemma}[Fractional moment balance]\label{lem:moment-balance}
For every solution in \cref{def:solution} and every $0<p<1$, $M_p$ is
locally absolutely continuous. Its weak balance is \eqref{eq:weak} with
$f(x)=x^p$, with absolutely integrable right-hand side on finite intervals.
\end{lemma}

\begin{proof}
H\"older's inequality gives $M_p(t)\leq N(t)^{1-p}m^p$.
Use the bounded Borel tests $f_R(x)=\min(x^p,R)$.
They are increasing, concave, and subadditive, with $f_R(0)=0$.
Consequently
\[
 0\leq f_R(x)+f_R(y)-f_R(x+y)\leq\min(x^p,y^p).
\]
After multiplication by $x+y$ this is bounded by $xy^p+yx^p$,
whose integral against $n_t\otimes n_t$ is $2mM_p(t)$.
For daughters, $f_R(z)/z$ is nonincreasing in $z$; since $z<x$ and
daughter mass is conserved,
\[
 0\leq\int_E f_R(z)b_{t,x}(\dd z)-f_R(x)
 \leq\int_E z^p b_{t,x}(\dd z)\leq2^{1-p}x^p.
\]
The last inequality is Jensen's inequality for the probability
$b_{t,x}/2$, whose mean is $x/2$. These bounds are uniform in $R$ and
integrable in time after multiplication by the locally integrable rates.
Monotone convergence on the left and dominated convergence on the right
of \eqref{eq:weak} yield the asserted identity. Its integrable right-hand
side also proves local absolute continuity.
\end{proof}

\subsection{A uniform estimate for controlled rates}

Write
\begin{equation}
 a_p=1+p-2^p>0\quad(0<p<1),\qquad
 \kappa=3-2\sqrt2=2a_{1/2}>0.
 \label{eq:constants}
\end{equation}
Positivity of $a_p$ follows from strict convexity of $2^p$ and its chord
between $p=0$ and $p=1$.

\begin{theorem}[Separation of the sampling laws]\label{thm:affinity}
Under \cref{ass:rates,def:solution}, for every $0<p<1$,
\begin{align}
 M_p'(t)&\leq
 \bigl[-b(t)(2-2^p)+\sigma(t)(2^{1-p}-1)\bigr]M_p(t)
 \quad\text{for almost every }t,\label{eq:raw-moment}\\
 \mathcal A_p(t)&\leq\mathcal A_p(0)
        e^{-a_p\Bc(t)-a_{1-p}F(t)}.\label{eq:affinity-decay}
\end{align}
In particular, writing $H(t)=M_{1/2}(t)$,
\begin{equation}
 \frac{H(t)^2}{N(t)}\leq\frac{H(0)^2}{N_0}e^{-\kappa C(t)}.
 \label{eq:half-decay}
\end{equation}
The coefficients are instantaneously sharp in the following sense: if
constants $c_1,c_2$ satisfy
$\mathcal A_p(t)\leq\mathcal A_p(0)e^{-c_1\Bc(t)-c_2F(t)}$ for all
$t\geq0$, all solutions of \cref{thm:existence}, and all daughter kernels
satisfying \eqref{eq:daughters}, then $c_1\leq a_p$ and $c_2\leq a_{1-p}$.
\end{theorem}

\begin{proof}
By \cref{lem:pair}, the coagulation term in the fractional moment balance
is at most $-\lambda(t)(2-2^p)mM_p(t)$. Jensen's inequality in the proof of
\cref{lem:moment-balance} bounds the fragmentation contribution by
$\sigma(t)(2^{1-p}-1)M_p(t)$. This proves \eqref{eq:raw-moment}.
Differentiate the number normalization using \eqref{eq:count}.
Its contribution is $-(1-p)(\sigma-b)$, and the resulting coefficients
are $-a_p b-a_{1-p}\sigma$. An integrating factor proves
\eqref{eq:affinity-decay} without dividing by $M_p$.
Taking $p=1/2$ and squaring gives \eqref{eq:half-decay}.

For the sharpness statement take constant rates, initial data
$n_0=N_0\delta_{x_0}$, and equal daughters $b_x=2\delta_{x/2}$, so that
$\mathcal A_p(0)=1$. For $f(x)=x^p$, weighted-variation continuity from
\cref{thm:existence}, the bound
$(x+y)|\Delta f(x,y)|\leq xy^p+yx^p\leq2(1+x)(1+y)$,
and the equal-split identity $\mathcal F f(x)=(2^{1-p}-1)x^p$ make the
moment-balance integrand continuous at zero. Thus $M_p$ has a right
derivative at zero, equal to the value of the vector field at $n_0$.
There both the pair inequality and Jensen's inequality are equalities, so
$M_p'(0+)=[-b(2-2^p)+\sigma(2^{1-p}-1)]M_p(0)$, and by \eqref{eq:count}
$\mathcal A_p'(0+)=-a_pb-a_{1-p}\sigma$. Now let $c_1,c_2$ satisfy the
stated bound. With $\sigma=0$ and $\lambda>0$ the bound gives
$\mathcal A_p(t)\leq1-c_1bt+o(t)$ as $t\downarrow0$, while
$\mathcal A_p(t)=1-a_pbt+o(t)$; hence $c_1\leq a_p$. With $\lambda=0$ and
$\sigma>0$ the same comparison gives $c_2\leq a_{1-p}$.

For equal splitting the identity $\int z^pb_x(\dd z)=2^{1-p}x^p$ holds for
every parent size $x$, not only at a monodisperse state. Hence with
$\lambda=0$, \cref{lem:moment-balance} gives
$M_p(t)=M_p(0)e^{(2^{1-p}-1)F(t)}$ and, by \eqref{eq:count},
$\mathcal A_p(t)=\mathcal A_p(0)e^{-a_{1-p}F(t)}$ for every initial measure
and every $t\geq0$. The fragmentation coefficient is then exact for all
times, not only instantaneously.
\end{proof}

\begin{remark}[What the estimate compares]\label{rem:relative}
If $C(t)\to\infty$, the sampling laws separate in total variation, because
\begin{equation}
 1-\norm{\eta_t-\pi_t}_{\TV}
 =\int_E\min\left(1,\frac{N(t)x}{m}\right)\eta_t(\dd x)
 \leq\mathcal A_p(t)\longrightarrow0.
 \label{eq:overlap-affinity}
\end{equation}
This concerns relative number and mass sampling. Outside the critical
case below it does not imply that raw mass moves to arbitrarily large
sizes.

Instantaneous sharpness does not identify the long-time exponent for a
prescribed daughter law or initial distribution, and the gap can be
large. For pure additive coagulation ($\sigma=0$, constant $\lambda$,
unit monodisperse initial data with $N_0=m=1$, so that $b=\lambda$) the
explicit solution of \citet{Golovin1963}, recorded in \citet{Bertoin2009},
has $N(t)=e^{-bt}$ and number law $\eta_t$ equal to the Borel law with
parameter $\mu_t=1-e^{-bt}$, that is
$\eta_t(\{k\})=e^{-\mu_tk}(\mu_tk)^{k-1}/k!$; for its probabilistic
interpretation see \citet{DeaconuTanre2000,Bertoin2002}, and for the
scaling theory of this kernel see \citet{MenonPego2008}. As $t\to\infty$
the number law converges to the Borel law with parameter one, whose tail
is of order $k^{-3/2}$. Consequently $M_p/N$ stays bounded for $p<1/2$ and
grows like $e^{(2p-1)bt}$ for $p>1/2$, and $\mathcal A_p=(M_p/N)N^p$
decays at the exact exponential rate $\min(p,1-p)\,b$, in the sense that
$-t^{-1}\log\mathcal A_p(t)\to\min(p,1-p)\,b$. For $p=1/2$ this is about
six times the coefficient $a_{1/2}\approx0.086$ in
\eqref{eq:affinity-decay}. The critical equal-split particle simulations
of \cref{sec:finite} show a gap of the same kind: for the largest
population the fitted late-time decay rate of $\mathcal A_{1/2}$ is about
$0.45b$, against the guaranteed $\kappa b\approx0.17b$. The estimate is
therefore a uniform Lyapunov bound, valid for every admissible daughter
law and initial measure, not an asymptotic rate.
\end{remark}

\subsection{Criticality: constant count and motion toward both boundaries}

Suppose now that $\lambda>0$ is constant and $\sigma=b$, so that both
rates are constant.
The daughter measure may still depend measurably on time and parent size.
Count remains exactly $N_0$. Set
\begin{equation}
 \kappa_p=3-2^p-2^{1-p}=a_p+a_{1-p}>0.
 \label{eq:kappap}
\end{equation}

\begin{corollary}[Two-boundary separation]\label{cor:critical}
At criticality, for every $0<p<1$,
\begin{equation}
 M_p(t)\leq M_p(0)e^{-b\kappa_p t}.
 \label{eq:critical-moment}
\end{equation}
For each fixed $p$ the coefficient is instantaneously sharp in the
sense of \cref{thm:affinity}: if $M_p(t)\leq M_p(0)e^{-cbt}$ for all
$t\geq0$, all critical solutions of \cref{thm:existence}, and all daughter
kernels satisfying \eqref{eq:daughters}, then $c\leq\kappa_p$. Among these
coefficients the unique maximum is $\kappa_{1/2}=\kappa$.
For every $\varepsilon,R>0$,
\begin{align}
 n_t([\varepsilon,\infty))
   &\leq\varepsilon^{-p}M_p(0)e^{-b\kappa_p t},\label{eq:number-window}\\
 \int_{(0,R]}x\,n_t(\dd x)
   &\leq R^{1-p}M_p(0)e^{-b\kappa_p t}.\label{eq:mass-window}
\end{align}
Viewed as measures on $[0,\infty)$, $n_t$ converges weakly to
$N_0\delta_0$. On the compactified space $[0,\infty]$, $\pi_t$ converges
weakly to $\delta_\infty$. No nonzero stationary solution with finite count
and positive finite mass exists.
\end{corollary}

\begin{proof}
Equation \eqref{eq:critical-moment} follows from \eqref{eq:raw-moment};
its sharpness follows from the monodisperse equal-split example in the
proof of \cref{thm:affinity}, where $M_p'(0+)=-b\kappa_pM_p(0)$, so
$M_p(t)=M_p(0)(1-b\kappa_pt+o(t))$ and $c\leq\kappa_p$.
The strict convexity and symmetry of $2^p+2^{1-p}$ give its unique minimum
$2\sqrt2$ at $p=1/2$, and its value is strictly below the common endpoint
value $3$ in $(0,1)$.
On $[\varepsilon,\infty)$, $x^p\geq\varepsilon^p$; on $(0,R]$,
$x\leq R^{1-p}x^p$. These prove the window bounds.
For a bounded continuous function on $[0,\infty)$, its oscillation near
zero is small and the number outside that neighborhood tends to zero.
This proves the first weak limit. A continuous function on
$[0,\infty]$ is uniformly close to its value at infinity outside a
sufficiently large bounded interval; \eqref{eq:mass-window} proves the
second limit. A stationary measure would have $M_p>0$ constant in time,
contradicting \eqref{eq:critical-moment}.
\end{proof}

Both limits are long-time statements. All material mass is conserved at
every finite time. The particles that become rare in number carry most
of the material. There is no assertion of finite-time shattering or
gelation. On a fixed bounded size interval with an upper cutoff $R$,
positive mass $m$ forces $M_p\geq m/R^{1-p}$; hence a discretization that
confines all mass there cannot satisfy \eqref{eq:critical-moment}
indefinitely.

\begin{corollary}[Exponentially moving observation windows]\label{cor:moving}
For $\varepsilon_0,R_0>0$ and $0<v<b\log2$, both
\begin{equation}
 n_t([\varepsilon_0e^{-vt},\infty))\longrightarrow0,\qquad
 \int_{(0,R_0e^{vt}]}x\,n_t(\dd x)\longrightarrow0
 \label{eq:moving}
\end{equation}
at exponential rates.
\end{corollary}

\begin{proof}
The two respective bounds, for arbitrary $p\in(0,1)$, are
\[
 \varepsilon_0^{-p}M_p(0)e^{-[b\kappa_p-pv]t},\qquad
 R_0^{1-p}M_p(0)e^{-[b\kappa_p-(1-p)v]t}.
\]
Since $\kappa_p/p\to\log2$ as $p\downarrow0$ and
$\kappa_p/(1-p)\to\log2$ as $p\uparrow1$, a suitable exponent exists
for each bound. They generally use different exponents $p$.
Since $\kappa_p$ is concave with $\kappa_0=0$, the quotient $\kappa_p/p$
decreases in $p$, so $\sup_{0<p<1}\kappa_p/p=\log2$; by the symmetry
$\kappa_p=\kappa_{1-p}$ the same holds for $\kappa_p/(1-p)$. Thus
$v<b\log2$ is exactly the range this argument reaches.
\end{proof}

These guaranteed velocities do not give a matching asymptotic front
speed. The endpoint $v=b\log2$ is not included.

\subsection{A nonstationarity criterion under rate perturbations}

\begin{proposition}[A sufficient open rate condition]\label{prop:robust}
Consider a jointly measurable symmetric kernel $(t,x,y)\mapsto K_t(x,y)$
and jointly measurable selection rate $(t,x)\mapsto S_t(x)$ satisfying
\begin{equation}
 a(x+y)\leq K_t(x,y)\leq A(x+y),\qquad
 0\leq S_t(x)\leq s_*,\qquad 0<a\leq A<\infty,
 \label{eq:robust-assumptions}
\end{equation}
with daughters as in \cref{ass:rates}. Use \cref{def:solution} with
$\lambda(s)(x+y)$ replaced by $K_s(x,y)$ in the coagulation term of
\eqref{eq:weak}, and $\sigma(s)\mathcal F_s f(x)$ replaced by
$S_s(x)\mathcal F_s f(x)$ in the fragmentation term. In particular, the
balance holds for every bounded Borel test, count is locally bounded,
and mass is conserved. For every such solution of mass $m$ and every
$0<p<1$,
\begin{equation}
 M_p(t)\leq M_p(0)e^{-\gamma_p t},\qquad
 \gamma_p=(2^{1-p}-1)(am\,2^p-s_*).
 \label{eq:robust}
\end{equation}
If $s_*<2am$, no stationary measure of mass $m$ has finite positive
particle count.
\end{proposition}

\begin{proof}
The count balance gives $N(t)\leq N_0e^{s_*t}$. The upper bounds in
\eqref{eq:robust-assumptions} and the proof of
\cref{lem:moment-balance} justify the fractional test. Its coagulation
contribution is at most $-am(2-2^p)M_p$ by the lower kernel bound,
and fragmentation contributes at most $s_*(2^{1-p}-1)M_p$.
Their difference equals $-\gamma_pM_p$ because
$2-2^p=2^p(2^{1-p}-1)$.
If $s_*<2am$, choose $p<1$ close enough to one that $am2^p>s_*$.
A stationary measure would then have a strictly positive moment that
decays exponentially, a contradiction.
\end{proof}

The threshold here is sufficient. We make no claim of its optimality or
of existence of stationary solutions above it. Count need not be constant
under \eqref{eq:robust-assumptions}.
